% I paragrafi di questo capitolo usano \paragraph* (forma "starred") invece di
% \paragraph: la forma non starred aggiungerebbe comunque una voce
% all'indice (il tocdepth=5 impostato in preface/5_table_of_contents.tex
% include anche il livello "paragraph"), cosa non voluta per queste voci
% puramente informative/redazionali. \paragraph* non aggiunge voci
% all'indice (né numera), esattamente come già fa \chapter* per il titolo
% del capitolo stesso poco sotto.
\cleardoublepage
\phantomsection
\pdfbookmark{Convenzioni tipografiche}{convenzioni-tipografiche}

\begingroup
\let\clearpage\relax
\let\cleardoublepage\relax

\chapter*{Convenzioni tipografiche}

Questa sezione riepiloga le convenzioni tipografiche adottate nel presente lavoro.

\paragraph*{Riferimenti a fonti giuridiche.} Le leggi, i decreti legislativi e le direttive europee citati nel testo (ad esempio la legge Stanca, il D.Lgs.~82/2022 o le direttive (UE) 2019/882 e (UE) 2016/2102) non seguono lo stile di citazione numerico adottato per le altre fonti, ma la convenzione propria della dottrina giuridica italiana: la nota a piè di pagina riporta per esteso \emph{tipo di atto, data completa e numero} (ed eventualmente l'articolo, il comma o l'allegato richiamati), ad esempio:
\begin{quote}
\itshape Legge 9 gennaio 2004, n.~4, art.~3, comma~1.
\end{quote}
Questa scelta rispecchia il modo in cui le fonti normative sono abitualmente richiamate nella letteratura giuridica e nei documenti ufficiali (AGID, Commissione europea), dove l'atto è identificato dagli estremi normativi e non da un numero progressivo di bibliografia.

\paragraph*{Riferimenti alle altre fonti (stile IEEE).} Tutte le fonti non giuridiche (standard tecnici come la norma EN~301549, documenti del W3C, linee guida AGID di natura tecnica, letteratura scientifica, rapporti come il WebAIM Million) sono invece citate secondo lo standard IEEE: la nota a piè di pagina riporta, tra parentesi quadre, il numero progressivo con cui la fonte compare in Bibliografia o in Sitografia, seguito dall'eventuale indicazione puntuale (capitolo, paragrafo, pagina), riportata invece fuori dalle parentesi, ad esempio:
\begin{quote}
\itshape [2], cap.~4, § 4.2.
\end{quote}
dove \emph{2} è il numero della fonte in bibliografia. Le fonti bibliografiche in senso stretto (testi, articoli) sono elencate in \emph{Bibliografia}; le pagine web e i documenti consultati online sono elencati separatamente in \emph{Sitografia}, con numerazione progressiva unica e condivisa fra le due sezioni.

\paragraph*{Termini stranieri.} I termini e le espressioni in lingua inglese privi di un equivalente italiano consolidato nel linguaggio tecnico (ad esempio \emph{screen reader}, \emph{backend}, \emph{prompt}, \emph{framework}) sono riportati in corsivo alla prima occorrenza e, per coerenza, in tutte le occorrenze successive. Non sono invece messi in corsivo gli acronimi (WCAG, LLM, API) né i termini ormai entrati nell'uso corrente italiano.

\paragraph*{Acronimi e abbreviazioni.} Gli acronimi e le abbreviazioni sono contrassegnati, alla loro comparsa nel testo, da un colore distintivo (ad esempio \textcolor{acronymColor}{WCAG}). Non tutti gli acronimi sono però accompagnati, nel testo, dalla propria forma estesa: quelli che compaiono in un contesto discorsivo sono di norma spiegati per esteso alla prima occorrenza (ad esempio \enquote{Web Content Accessibility Guidelines (WCAG)}), e in forma abbreviata nelle occorrenze successive; quelli che compaiono invece all'interno di una denominazione ufficiale (il nome di una norma tecnica, di un atto normativo o di un ente) restano sempre nella sola forma abbreviata, poiché espanderli altererebbe la denominazione stessa (ad esempio \enquote{norma \textcolor{acronymColor}{UNI} \textcolor{acronymColor}{CEI}~EN~301549}, e non \enquote{norma Ente Italiano di Normazione Comitato Elettrotecnico Italiano EN~301549}): in questo caso la forma estesa non compare nel testo, ma è comunque reperibile, tramite il collegamento ipertestuale associato alla sigla, nell'elenco degli Acronimi e abbreviazioni in fondo al documento.
 

\paragraph*{Termini del glossario.} I termini tecnici che ricevono una definizione puntuale nel Glossario sono contrassegnati, alla loro comparsa nel testo, da un apice \textsubscript{\textbf{\textit{G}}} (ad esempio \emph{termine}\glox).

\paragraph*{Citazioni testuali.} Le citazioni testuali tratte da fonti normative o tecniche sono racchiuse fra virgolette basse, o caporali (\enquote{esempio}), per distinguerle dal corsivo usato per i termini stranieri e dall'apice usato per i rimandi al glossario.

\paragraph*{Adattamenti all'interno di una citazione testuale.} Quando, all'interno di una citazione testuale, è necessario modificare la coniugazione di un verbo (o un altro elemento grammaticale) per adattare la citazione alla struttura sintattica della frase che la contiene, la porzione modificata rispetto al testo originale della fonte è racchiusa fra parentesi quadre, secondo la convenzione filologica standard, ad esempio:
\begin{quote}
\itshape \enquote{i componenti dell'interfaccia utente e la navigazione [sono] operabili}
\end{quote}
Le parentesi quadre segnalano quindi un intervento del redattore (ad esempio un cambio di persona, di numero o di tempo verbale) rispetto al testo originale, senza alterarne il significato.

\paragraph*{Codice sorgente.} Gli estratti di codice sono composti in carattere monospazio con evidenziazione sintattica, numerati progressivamente per capitolo ed elencati nel relativo Elenco dei codici sorgenti.
\endgroup